%% kapa: exam=1 round=33 year=2022 period=2 form=A subject=appraisal_law subject_ko=감정평가관계법규 seq=1 number=1
%% source: Q-Net 감정평가사 기출문제 (cst00309) artlSeq=5209997 file=raw/1008/5209997/2022년도 제33회 감정평가사 1차시험 2교시(감정평가 관계법규, 회계학).pdf
%% conversion: kordoc 4.24.51 -> tools/gen1.py (auto); review=proofread-sonnet; flags=none
\begin{question}{1}
국토의 계획 및 이용에 관한 법령상 국토교통부장관이 도시·군관리계획의 수립 기준을 정할 때 고려하여야 하는 사항이 아닌 것은?
\begin{choices}
  \item 공간구조는 생활권 단위로 적정하게 구분하고 생활권별로 생활·편익시설이 고루 갖추어지도록 할 것
  \item 녹지축·생태계·산림·경관 등 양호한 자연환경과 우량농지, 문화재 및 역사문화환경 등을 고려하여 토지이용계획을 수립하도록 할 것
  \item 수도권안의 인구집중유발시설이 수도권외의 지역으로 이전하는 경우 종전의 대지에 대하여는 그 시설의 지방이전이 촉진될 수 있도록 토지이용계획을 수립하도록 할 것
  \item 도시의 개발 또는 기반시설의 설치 등이 환경에 미치는 영향을 미리 검토하는 등 계획과 환경의 유기적 연관성을 높여 건전하고 지속가능한 도시발전을 도모하도록 할 것
  \item 광역계획권의 미래상과 이를 실현할 수 있는 체계화된 전략을 제시하고 국토종합계획 등과 서로 연계되도록 할 것
\end{choices}
\end{question}
\answer{5}
